% =========================================================================
% SECTION 4: ENGINEERING EVALUATION
% =========================================================================
\section{Engineering Evaluation}
\label{sec:evaluation}

The engineering evaluation investigates whether the implemented PRISM Android application satisfies its kinematic, geometric, and multimodal warning specifications under controlled indoor test conditions.

\subsection{Evaluation Environment and Hardware Platform}
Testing was conducted on commodity Android smartphones running Android 14 (API 34, minimum SDK 26) equipped with autofocus rear cameras and 6-axis IMUs (accelerometer and gyroscope). The physical testbed comprised structured indoor institutional corridors ($2.4\text{--}3.0$\,m hallway width) with vinyl tile flooring, standard overhead fluorescent lighting, and acoustic reverberation representative of public buildings.

% Table 2: Controlled Indoor Evaluation Scenarios and Dataset
\begin{table}[t]
\caption{Controlled Indoor Engineering Evaluation Scenarios and Dataset}
\label{tab:evaluation_setup}
\centering
\scriptsize
\setlength{\tabcolsep}{3pt}
\begin{tabularx}{\columnwidth}{p{1.9cm}ccp{3.7cm}}
\toprule
\textbf{Scenario Family} & \textbf{Runs} & \textbf{Frames} & \textbf{Kinematic Profile / Obstacle Configuration} \\
\midrule
Stationary Obstacles & 3 & 630 & Static box / mannequin placed at $0.5, 1.0, 2.0, 3.0$\,m \\
Head-On Approach & 3 & 411 & Pedestrian approaching user collinear at $1.0\text{--}1.4$\,m/s \\
Lateral Crossing & 4 & 561 & Pedestrian crossing perpendicularly ($2$ collision, $2$ miss) \\
Moving Away & 2 & 385 & Co-directional walking; tests false-warning suppression \\
Direction Change & 2 & 283 & Pedestrian reverses heading; tests Kalman coasting \\
Multiple Pedestrians & 3 & 532 & $\ge 2$ simultaneous agents; tests threat prioritization \\
\midrule
\textbf{Total Video Runs} & \textbf{17} & \textbf{2,827} & \textbf{Evaluated across real indoor corridor scenarios} \\
\midrule
A/B Mode Evaluation & 10 & --- & Controlled stationary, crossing, and evasion trials \\
\bottomrule
\end{tabularx}
\vspace{-0.15cm}
\par{\raggedright \tiny Evaluated on Android device (API 34) under controlled indoor engineering conditions.\par}
\vspace{-0.3cm}
\end{table}


\subsection{Software Test Suite Verification}
System correctness and regression safety were verified via the automated test suite in \texttt{PrismBridgeUnitTest.kt}. The suite contains 100 dedicated unit and integration tests executing directly on the Android runtime:
\begin{itemize}
    \item \textbf{Tests 1--62}: Core ground-plane projection geometry, Mahalanobis gating, multi-target tracking, TTC/CPA calculations, and collision decision arbitration.
    \item \textbf{Tests 63--71}: IMU pitch/roll orientation fusion, gravity filtering, and dynamic projection tilt compensation.
    \item \textbf{Tests 72--82}: Nine-state hazard lifecycle machine, escalation, user-stopping detection, and hold persistence.
    \item \textbf{Tests 83--89}: Accelerometer step detection, refractory timing, and 6-state user motion lifecycle modeling.
    \item \textbf{Tests 90--100}: Spatial audio stereo panning, dual-harmonic emergency synthesis, and multisensory mode arbitration.
\end{itemize}

\subsection{Controlled Scenario Dataset and Real-World Video}
The system was evaluated across 17 recorded real-world indoor video sequences totaling 2,827 frames ($720 \times 1280$ at 30\,FPS) spanning six canonical pedestrian encounter scenarios (Table~\ref{tab:evaluation_setup}): stationary obstacles at varying ranges ($0.5\text{--}3.0$\,m), collinear head-on approaches ($1.0\text{--}1.4$\,m/s closing speed), lateral crossings, departing entities, sudden trajectory reversals, and multi-obstacle scenes. Ground-truth obstacle positions and collision breach timestamps were established via synchronized rangefinder measurement and frame-level video annotation.

\subsection{Multisensory A/B Evaluation Protocol}
To evaluate the functional utility of stereo spatial audio, a controlled A/B engineering comparison was conducted between \texttt{MODE\_A\_SPEECH\_HAPTIC} and \texttt{MODE\_B\_SPEECH\_SPATIAL\_HAPTIC} across 10 controlled obstacle encounters. Sighted engineering evaluators navigated a standard corridor while wearing stereo over-ear headphones. Six evaluation metrics were collected:
\begin{enumerate}
    \item \textbf{Direction Identification Accuracy (\%)}: Percentage of trials where the evaluator stepped in the system-recommended avoidance path.
    \item \textbf{Wrong-Direction Rate (\%)}: Percentage of trials where the initial steering response was opposite to the clear path.
    \item \textbf{Reaction Time (ms)}: Interval from sensory stimulus onset to first evasive movement, measured via synchronized video frame analysis.
    \item \textbf{Warning Comprehension (\%)}: Accurate verbal identification of the encounter hazard type post-trial.
    \item \textbf{NASA-TLX Score (0--100)}: 6-subscale subjective cognitive workload assessment (lower indicates reduced cognitive burden).
    \item \textbf{SUS Score (0--100)}: Standard 10-item System Usability Scale instrument (scores $\ge 68$ denote above-average usability).
\end{enumerate}

\textbf{Participant Distinction}: This evaluation was conducted with engineering evaluators to verify sensory pipeline behavior. It does not constitute a formal human-subject clinical trial with BVI participants.
